\documentclass[12pt,letterpaper]{report}
\usepackage{graphicx}
\usepackage{scrextend}
\usepackage{vmargin}
\usepackage{graphicx}
\usepackage{multirow}
\usepackage[utf8]{inputenc}
\usepackage[spanish]{babel}
\usepackage{multicol}
\usepackage{enumerate}
\usepackage{float}
\usepackage{amsmath, amsthm, amssymb, amsfonts}
\usepackage[usenames]{color}
\parindent=0mm
\pagestyle{empty}
\definecolor{miorange}{rgb}{0.91, 0.43, 0.0}
\begin{document}
\setmargins{2.5cm}      
{1.5cm}                     
{2cm}  
{24cm}                    
{10pt}                          
{1cm}                          
{0pt}                             
{2cm}
\begin{flushright}
\today
\end{flushright}
\vspace{-1.75cm}
\section*{Tarea Clase 6}
\subsection*{Obtenga las soluciones de las siguientes ecuaciones:}
\begin{enumerate}
\item \begin{equation*}
    2x-(5x+3)=7+(3x-2)
\end{equation*}
\item \begin{equation*}
    5x-9=3(x-2)
\end{equation*}
\item \begin{equation*}
    -2+x-8=2-x+8
\end{equation*}
\item \begin{equation*}
    2x+x+5=18-3x
\end{equation*}
\item \begin{equation*}
    -2x-5=8x+3
\end{equation*}
\item \begin{equation*}
    2x+3(x-1)=6-4(2x+3)
\end{equation*}
\item \begin{equation*}
    3x+5=\frac{5}{2}x+4
\end{equation*}
\item \begin{equation*}
    4(y+1)+3y-1=7y+3
\end{equation*}
\item \begin{equation*}
    8x-15x-30x-51x=53x+31x-172
\end{equation*}
\item \begin{equation*}
    \frac{3}{2}x+\frac{1}{4}+2=\frac{3}{4}x-\frac{1}{3}x
\end{equation*}
\item \begin{equation*}
    x^2+6=7x
\end{equation*}
\item \begin{equation*}
    3x^2+2x-8=0
\end{equation*}
\item \begin{equation*}
    2x^2-12x+18=0
\end{equation*}
\item \begin{equation*}
    -x^2=-3x-10
\end{equation*}
\item \begin{equation*}
    x^2-2x=1
\end{equation*}
\item \begin{equation*}
    -3x^2-6x+9=0
\end{equation*}
\item \begin{equation*}
    -2x^2-10x=0
\end{equation*}
\item \begin{equation*}
    3x^2=-6x-3
\end{equation*}
\end{enumerate}
Las respuestas de las tareas pueden ser enviadas al siguiente correo giovannilopez9808@gmail.com con el asunto entre corchetes de la siguiente manera [Tarea \#] y el nombre del archivo con sus apellidos y nombre, ya sea el nombre completo o solo una parte, por ejemplo \textit{LopezPadilla.pdf}, \textit{LopezGiovanni.pdf}, \textit{LopezPadillaGiovanniGamaliel.pdf}. Las respuestas envienlas en formato pdf, estas pueden ser totalmente digitales (que escriban todo en la computadora) o que sean fotografías, pero por favor, envielo en solo un documento.
\end{document}